\documentclass[tikz,border=2pt]{standalone}
\usetikzlibrary{graphs,graphdrawing,arrows.meta,shapes.geometric}
\usegdlibrary{layered}
\tikzset{
  tarefa/.style={draw, rounded corners=2pt, align=center, text width=1.7cm,
                 minimum height=0.8cm, font=\scriptsize},
  inicio/.style={circle, draw, minimum size=0.45cm, inner sep=0pt},
  fim/.style={circle, draw, very thick, minimum size=0.45cm, inner sep=0pt},
  meio/.style={circle, draw, double, minimum size=0.45cm, inner sep=0pt},
  decisao/.style={diamond, draw, minimum size=0.6cm, inner sep=0pt},
  rotulo/.style={font=\tiny\itshape, align=center},
  condicao/.style={font=\tiny\itshape, text=black, fill=white, inner sep=0.5pt},
}
\begin{document}
\begin{tikzpicture}[>={Latex[length=1.4mm]}]
\graph[layered layout, grow=right, level distance=2.0cm, sibling distance=1.3cm, edges={gray!70, ->}] {
  n1[inicio, as={}];
  n2[tarefa, as={File complaint}];
  n3[tarefa, as={Log complaint and\\assign to department}];
  n4[tarefa, as={Review details}];
  n5[decisao, as={$\times$}, label={[rotulo]above:Refund justified?}];
  n6[tarefa, as={Notify customer of\\decision}];
  n7[tarefa, as={Process reimbursement}];
  n8[fim, as={}];
  n9[fim, as={}];
  n10[tarefa, as={Provide feedback}];
  n11[fim, as={}];
  n1 -> n2;
  n2 -> n3;
  n3 -> n4;
  n4 -> n5;
  n5 ->[edge node={node[condicao] {yes}}] n6;
  n5 ->[edge node={node[condicao] {no}}, bend right=35] n6;
  n6 -> n7;
  n7 -> n8;
  n6 -> n9;
  n8 -> n10;
  n9 -> n10;
  n10 -> n11;
};
\end{tikzpicture}
\end{document}
